% Five-slide Black--Scholes validation block: Kaiden Wessels, 6:50.

\speakerlabel{Kaiden}{1:10}
\begin{frame}{Black--Scholes is the validation gate}
  \academicbox[\textwidth]{\centering
    {\color{DeepNavy}\small\bfseries Black--Scholes is the controlled validation gate: the correct finite-grid benchmark is known.}}
  \vspace{0.12cm}
  \begin{columns}[T,onlytextwidth]
    \begin{column}{0.58\textwidth}
      \begin{beamercolorbox}[rounded=true,sep=0.20cm]{block body}
        {\fontsize{8.8}{10.5}\selectfont
        \[
          \begin{aligned}
          \he
          &=
          \pi+
          \underbrace{\sum_{n=0}^{N-1}\delta_n
          \bigl(S_{t_{n+1}}-S_{t_n}\bigr)}_{\text{self-financing gains}}
          -(S_T-K)^+,\\[-0.04cm]
          &\hspace{1.15cm}\min_{\pi,\theta}\;\mathbb{E}[\he^2].
          \end{aligned}
        \]}
      \end{beamercolorbox}
      \vspace{0.12cm}
      \claimbox{The analytic delta is never used as a training target.}
    \end{column}
    \begin{column}{0.38\textwidth}
      {\fontsize{8.5}{10.3}\selectfont
      \begin{tightitemize}
        \item \(\delta_n=f_\theta(x_n)\) is adapted to information at hedge date \(t_n\).
        \item Negative \(\he\) is a seller shortfall; seller loss is \(L=-\he\).
        \item With \(r=0\), gains are undiscounted in the submitted experiment.
      \end{tightitemize}}
    \end{column}
  \end{columns}
  \vspace{0.12cm}
  \begin{beamercolorbox}[rounded=true,sep=0.11cm]{block body}
    \centering{\small
      \(S_0=1\quad K=0.9\quad \sigma=0.4\quad T=0.5\quad r=0\quad N=125\)}
  \end{beamercolorbox}
  \sourcefoot{Submitted report, Secs.\ 1.1, 2.2.2 and 3.1; Eqns.\ (1.1)--(1.2).}
  \note{Frame Black--Scholes as the validation gate, then define premium, position, self-financing gains and payoff. Emphasise that MSE is the submitted terminal objective and that no analytic delta labels enter training.}
\end{frame}

\speakerlabel{Kaiden}{1:25}
\begin{frame}{Training is pathwise and differentiable end to end}
  \centering
  \begin{tikzpicture}[
    node distance=0.17cm,
    box/.style={rounded corners=3pt,draw=DeepNavy!20,fill=Mist,
      minimum height=0.82cm,text width=2.23cm,align=center,
      font=\fontsize{7.6}{8.7}\selectfont\bfseries},
    arrow/.style={-{Latex[length=2mm]},line width=0.9pt,color=Teal}
  ]
    \node[box] (sim) {simulate path};
    \node[box,right=of sim] (state) {form causal state};
    \node[box,right=of state] (mlp) {shared MLP};
    \node[box,right=of mlp] (pos) {positions};
    \node[box,below=0.38cm of pos] (gain) {self-financing gains};
    \node[box,left=of gain] (err) {terminal error};
    \node[box,left=of err] (mse) {MSE};
    \node[box,left=of mse] (bp) {backpropagation};
    \draw[arrow] (sim)--(state);
    \draw[arrow] (state)--(mlp);
    \draw[arrow] (mlp)--(pos);
    \draw[arrow] (pos)--(gain);
    \draw[arrow] (gain)--(err);
    \draw[arrow] (err)--(mse);
    \draw[arrow] (mse)--(bp);
    \draw[arrow] (bp.west) -- ++(-0.32,0) |- (mlp.south);
  \end{tikzpicture}

  \vspace{0.15cm}
  \begin{columns}[T,onlytextwidth]
    \begin{column}{0.48\textwidth}
      \academicbox{\centering\small\bfseries One sample = one complete trajectory}
    \end{column}
    \begin{column}{0.48\textwidth}
      \academicbox{\centering\small Black--Scholes paths use the exact lognormal grid transition.}
    \end{column}
  \end{columns}

  \vspace{0.12cm}
  \begin{beamercolorbox}[rounded=true,sep=0.12cm]{block body}
    \centering{\small\bfseries
      100,000 training paths\quad\textcolor{Teal}{|}\quad
      30,000 validation paths\quad\textcolor{Teal}{|}\quad
      100,000 held-out test paths}
  \end{beamercolorbox}
  \vspace{0.10cm}
  \claimbox{Supervision is financial: payoff minus self-financing wealth.}
  \sourcefoot{Submitted report, Algorithm A.6.1, Appendix A.1 and Listings C.1--C.3.}
  \note{Walk through one complete trajectory. Validation loss selects the retained epoch; the final numbers are then computed on freshly generated held-out test paths. Mention that the exact lognormal transition removes time-discretisation error from Black--Scholes path generation.}
\end{frame}

\speakerlabel{Kaiden}{1:15}
\begin{frame}{Architecture selection is validation-based}
  \begin{columns}[T,onlytextwidth]
    \begin{column}{0.54\textwidth}
      \begin{tikzpicture}[
        layer/.style={rounded corners=3pt,draw=DeepNavy!20,fill=Mist,
          minimum width=1.35cm,minimum height=0.88cm,align=center,font=\scriptsize},
        arr/.style={-{Latex[length=2mm]},color=Teal,line width=0.9pt},
        node distance=0.22cm
      ]
        \node[layer] (x) {\(\log(S/K)\)\\\(\tau/T\)};
        \node[layer,right=of x] (l1) {64\\tanh};
        \node[layer,right=of l1] (l2) {64\\tanh};
        \node[layer,right=of l2] (l3) {64\\tanh};
        \node[layer,right=of l3,fill=Teal!12] (o) {1\\sigmoid};
        \draw[arr] (x)--(l1);\draw[arr] (l1)--(l2);\draw[arr] (l2)--(l3);\draw[arr] (l3)--(o);
        \node[below=0.24cm of l2,font=\scriptsize,color=Slate] {one shared map at all 125 hedge dates};
      \end{tikzpicture}
      \vspace{0.27cm}
      {\fontsize{8.3}{9.9}\selectfont
      \begin{tightitemize}
        \item \textbf{Log-moneyness} is scale-normalised; \textbf{time-to-maturity} is explicit.
        \item \textbf{Shared weights} reflect the Markov, time-homogeneous structure.
        \item \textbf{Sigmoid} enforces the economically justified call-delta range \(0<\delta<1\).
        \item A scalar premium is learned jointly with the hedge.
      \end{tightitemize}}
    \end{column}
    \begin{column}{0.42\textwidth}
      \academicbox{\centering
        {\color{Teal}\fontsize{15}{17}\selectfont\bfseries 12 candidates \(\times\) 3 seeds\par}
        \vspace{0.05cm}
        {\small Selected on mean validation loss}}
      \vspace{0.12cm}
      \academicbox{\centering\small\bfseries Independently retrained for the final benchmark}
      \vspace{0.12cm}
      {\centering\scriptsize Illustrative close candidates\par}
      \vspace{0.03cm}
      \begin{table}
        \centering
        \fontsize{6.7}{7.7}\selectfont
        \setlength{\tabcolsep}{2.4pt}
        \begin{tabular}{lr}
          \toprule
          Candidate & Mean test RMSE\\
          \midrule
          \rowcolor{Teal!10} Shared norm.\ \(64{\times}3\), tanh & 0.007891\\
          Shared norm.\ \(64{\times}2\), ReLU & 0.007912\\
          Time-separated \(16{\times}1\) & 0.008083\\
          \bottomrule
        \end{tabular}
      \end{table}
      \vspace{0.04cm}
      {\scriptsize\color{Coral}\bfseries The selected model was not uniquely dominant.\par}
    \end{column}
  \end{columns}
  \sourcefoot{Submitted report, Secs.\ 3.1.2--3.1.3; Tables 3.1 and A.3--A.6.}
  \note{Explain the economic constraints first. Then state that all 12 candidates were compared over three independent seeds using mean validation loss, and that the selected specification was independently retrained. The table is only a visual reminder that close alternatives remained.}
\end{frame}

\speakerlabel{Kaiden}{1:20}
\begin{frame}{The fair benchmark matches the grid and objective}
  \begin{columns}[T,onlytextwidth]
    \begin{column}{0.58\textwidth}
      \begin{tikzpicture}[
        item/.style={rounded corners=3pt,minimum height=0.68cm,text width=6.35cm,
          align=left,font=\small,inner xsep=7pt},
        node distance=0.10cm
      ]
        \node[item,fill=Coral!10,text=Ink] (u) {\textbf{No hedge}\\[-0.04cm]{\scriptsize raw option risk}};
        \node[item,fill=Gold!12,below=of u] (bs) {\textbf{Black--Scholes delta}\\[-0.04cm]{\scriptsize continuous-time rule sampled on the grid}};
        \node[item,fill=Teal!18,below=of bs] (dt) {\textbf{Discrete-time MSE-optimal}\\[-0.04cm]{\scriptsize conditional projection on the same grid}};
        \node[item,fill=Mist,below=of dt] (poly) {\textbf{Degree-2 polynomial}\\[-0.04cm]{\scriptsize low-capacity approximation}};
        \node[item,fill=DeepNavy!12,below=of poly] (nn) {\textbf{Neural hedge}\\[-0.04cm]{\scriptsize terminal-MSE training without delta labels}};
      \end{tikzpicture}
    \end{column}
    \begin{column}{0.38\textwidth}
      \academicbox{\centering
        {\small\bfseries Discrete-time conditional projection\par}
        \vspace{0.05cm}
        {\fontsize{7.5}{8.7}\selectfont\(
          \displaystyle\deltadt_n=
          \frac{\operatorname{Cov}\!\left(g(S_T),\Delta S_n\mid S_{t_n}\right)}
          {\operatorname{Var}\!\left(\Delta S_n\mid S_{t_n}\right)}
        \)}}
      \vspace{0.12cm}
      \academicbox{\centering
        {\color{Teal}\small\bfseries Matches both the trading grid and the terminal-MSE objective}}
    \end{column}
  \end{columns}
  \vspace{0.16cm}
  \claimbox{The comparison asks whether the network recovers the finite-grid solution---not whether it beats correct Black--Scholes theory.}
  \sourcefoot{Submitted report, Sec.\ 3.1.1 and Appendix A.4; Schweizer (1995, 2001).}
  \note{Make the discrete-time conditional projection visually and conceptually central. Black--Scholes delta is a continuous-time rule sampled on the grid; the discrete-time hedge matches both the submitted grid and terminal-MSE objective. The next slide gives the held-out result.}
\end{frame}
